% Draft for the next numbered statement; publication source is frozen.
% This proves the actual framed finite-field group and the three possible
% Kummer comparisons.  It does not choose an arithmetic eta class.

\paragraph{The actual framed point group.}
Put $k=\mathbb F_4=\mathbb F_2[\zeta]$,
$E_0:y^2+y=x^3$, and let $D=E_0(k)=E_0[3]$.
The coordinate change and ordered cusps are those of B.35.
The inverse change is
\[
                  [X:Y:Z]=[y+\zeta:y+\zeta^2:x].
\]
Write $O$ for infinity.  For every $P\in D$ let $z_P$
be the point of the generalized Jacobian represented by
$\mathcal O(P-O)$ with the prescribed frames
$s_P^{-1}$ at $P$, $s_O$ at $O$, and one elsewhere;
put $z_O=0$.
The split torus is $(k^\times)^D/k^\times$, whose
coordinates, normalized at $O$, are written as
$W_0=\mathbb F_3^8$ using $\zeta$.
For $R\in D$ write $\epsilon_R(a)$ for the torus vector
having exponent $a$ at $R$ and zero elsewhere, modulo
the diagonal; in particular
$\epsilon_O(a)=(-a,\ldots,-a)$.
Every generalized-Jacobian point has a unique expression
$z_P+t$, with $P\in D$ and $t\in W_0$.
Indeed its compact image determines $P$, and the kernel
is the displayed split torus.  Surjectivity onto $E_0(k)$
is also explicit from the framed representatives $z_P$.

The point group is obtained without any choice of a
Tate division.  For $P,Q\ne O$ put $R=P+Q$.
If $R\ne O$, let
\[
 f_{P,Q}=\ell_{P,Q}/(x-x_R),
\]
where $\ell_{P,Q}$ is the secant or tangent line with
coefficient one on $y$.
If $R=O$, put $f_{P,Q}=x-x_P$; if either argument is
$O$, put $f_{P,Q}=1$.
In every case
\[
          \operatorname{div}(f_{P,Q})=P+Q-R-O.
\]
Multiplication by this function trivializes the
corresponding principal line bundle.  It sends its
prescribed frame at a cusp $T$ to the leading value
of $f_{P,Q}$ in $s_T$.  Therefore
\[
 z_P+z_Q-z_R=c(P,Q),\qquad
 c(P,Q)_T=
 \log_\zeta\frac{\operatorname{lead}_{s_T}f_{P,Q}}
                         {\operatorname{lead}_{s_O}f_{P,Q}}.
                                                               \tag{TP.1}
\]
This proves the direction and the sign of the toric
correction, rather than inferring it from an abstract
group presentation.

Here are complete finite formulas for these leading
values.  At a finite cusp $T=(x_T,y_T)$ set
$\delta=x-x_T$.  The curve equation gives
\[
 y-y_T=x_T^2\delta+\delta^3+\delta^6+\cdots.
                                                               \tag{TP.2}
\]
For a nonvertical line of slope $m$, its order and
leading value are read first from its constant term,
then from $x_T^2+m$, and finally from its cubic
coefficient one.  A vertical line has its ordinary
constant term or leading term $\delta$.
The actual parameters from the original plane model
are
\[
 s_{A_i}=(y+\zeta)/x,\quad
 s_{B_i}=(y+\zeta^2)/x,\quad
 s_{C_i}=x/(y+\zeta).
\]
Their derivatives relative to $dx$ are
\[
 k_T=\frac{ds_T}{dx}(T)=
 \begin{cases}
 x_T,&T=A_i\text{ or }B_i,\\
 (y_T+\zeta)^{-1},&T=C_1,C_2,\\
 1,&T=O.
 \end{cases}                                                   \tag{TP.3}
\]
For the last equality, put $q=x/(y+\zeta)$ and
substitute $y=x/q+\zeta$ into the cubic equation.
It gives $x=q^{-2}+q+O(q^7)$, so $dx/dq$ has leading
coefficient one.  Also $y$ has leading term $q^{-3}$.
Thus all the functions in (TP.1) have leading coefficient
one at $O$.
If their order at $T$ is $r$, multiply their leading
coefficient in $\delta$ by $k_T^{-r}$ to obtain their
leading coefficient in the prescribed parameter.

The program \path{c_char2_translation_paths.py} performs
all these operations in $k$ itself.
It checks the cocycle identity for $c$ on all $729$
triples and its symmetry on all $81$ pairs.
Consequently
\[
 (P,a)+(Q,b)=(P+Q,a+b+c(P,Q))                                \tag{TP.4}
\]
is the complete framed group law, with $3^{10}$ elements.
For the compact basis $P_1=A_0$, $P_2=A_2$, it verifies
\[
                    3z_P=2U\,\operatorname{coord}(P)
                         \quad\text{in }W_0.                \tag{TP.5}
\]
The two columns of $U$ are independent modulo three.
Hence $3J_{E,D}(k)$ has rank two over $\mathbb F_3$,
and (TP.4) has type $C_9^2\oplus C_3^6$.

\paragraph{Automorphisms and their tangent paths.}
For a geometric automorphism $g$ in the subgroup
$E_0[3]\rtimes(Q_8\rtimes C_3)$ of order $216$, write
$g^*dx=c_gdx$.  The pushed local parameter at $gP$
is $s_P\circ g^{-1}$, whose derivative relative to
$s_{gP}$ is $k_P/(c_gk_{gP})$.
Thus the exact pushforward of the framed point is
\[
 \begin{split}
 g(z_P)={}&z_{gP}-z_{gO}
       +\epsilon_{gP}\left(\log_\zeta
                              \frac{c_gk_{gP}}{k_P}\right)\\
       &+\epsilon_{gO}\left(\log_\zeta
                              \frac1{c_gk_{gO}}\right).
 \end{split}                                                   \tag{TP.6}
\]
On the torus $g$ permutes the nine coordinates before
renormalizing the $O$ coordinate.
Formula (TP.6) is verified on all group laws for the
eight quaternion elements, the order-three generator,
and the two generating translations: $891$ checks.
It supplies genuine framed automorphism data, not
only cusp permutations.

For translation by $d$, $c_g=1$ and the image of the
original unit tangent has derivative $k_d$ at $d$.
The path cocycle in abelianization therefore has the
actual local Kummer class of
\[
                         z_d+\epsilon_d(\log_\zeta k_d).
                                                               \tag{TP.7}
\]
For an origin-fixing automorphism this formula becomes
the toric class of its tangent multiplier.  In particular
it is zero for $Q_8$ and is
$\kappa_3(\zeta)e_O$ for $\sigma$, agreeing with
the independently proved tangent-rescaling identity.

Arithmetic Frobenius two is $(x,y)\mapsto(x^2,y^2+1)$.
The original plane parameters are defined over
$\mathbb F_2$, so it sends $z_P$ to $z_{F_2P}$ with
no additional frame multiplier; on the torus it squares
the values and permutes the cusps.
All $81$ corresponding group-law identities are checked.

\paragraph{The remaining three Kummer comparisons.}
Let $H=V/(F_4-I)V$ in the fixed integral lattice, with
\[
           F_4-I=\begin{pmatrix}-3I_2&0\\-2U&3I_8\end{pmatrix}.
                                                               \tag{TP.8}
\]
A pair of integer vectors $(a,b)$ has unique normal form
\[
       (a\bmod3,\ b-2U\lfloor a/3\rfloor\bmod3)
                                                               \tag{TP.9}
\]
in this finite group.
The canonical local Kummer isomorphism
$J_{E,D}(k)\xrightarrow{\sim}H$ sends a torus value
$\zeta$ to its corresponding peripheral coordinate.
On the compact quotient $F_4-I=-3$, so it sends
$P_i$ to $-f_i$ modulo the boundary; this sign is
forced by the Kummer equation $FQ-Q=P$.

Every comparison respecting these toric and compact
maps is therefore determined by
\[
                 z_{P_i}\longmapsto-f_i+b_i,
                         \qquad b_i\in\mathbb F_3^8.
                                                               \tag{TP.10}
\]
Conversely (TP.5) and (TP.8) show that every such pair
defines a group isomorphism with the required graded
maps.  Equivariance under each actual automorphism
and Frobenius two imposes linear equations on the
sixteen entries of $b_1,b_2$.
The script constructs these equations from
(TP.4)--(TP.10), using the compatible quaternion,
CM, and Frobenius matrices of B.35.
It retains an integral $-1$ in every boundary column
which maps a cusp to $O$ before dividing the cross
block by three; replacing it early by the residue
$2$ would change the equation.

Exact elimination over $\mathbb F_3$ gives rank fifteen,
with particular solution
\[
 b^{(0)}=(1,2,2,2,1,2,1,1;\ 0,2,0,0,2,2,2,0)
\]
and kernel generator
\[
 b^{(1)}=(0,1,2,0,2,1,2,1;\ 1,2,0,2,1,0,2,1).
                                                               \tag{TP.11}
\]
Thus the actual comparison is one of precisely
$b^{(0)}+a b^{(1)}$, $a\in\mathbb F_3$.
For each of them (TP.7) gives an explicit actual-path
candidate in the fixed coefficient module.
The assertion is an exhaustive set of three possible
comparisons, not a claim that one arbitrarily selected
solution is the canonical Kummer map.
Any ensuing affine test of $\eta$ must retain all three,
unless an additional division normalization selects one.
